\section{Discussion}

\subsection{What Truco Reveals Beyond Poker}

% Escalation as a test of meta-game reasoning
% Discrete vs continuous bluffing
% Compact game tree enables more comprehensive analysis

\subsection{Escalation as Meta-Reasoning}

% The escalation chain tests a specific form of reasoning: commitment under uncertainty
% Models must reason about credibility, counter-escalation probability, and expected value
% This is distinct from poker's betting rounds

\subsection{Cultural Priors in Model Responses}

% Do models trained primarily on English data perform differently on a Brazilian game?
% Language bias results discussion

\subsection{Limitations}

Our work has several limitations:
\begin{itemize}
    \item \textbf{Two-player only.} Traditional Truco Paulista is often played as a 4-player team game with partner signaling (sinais). The team variant introduces cooperative reasoning that is absent from our evaluation.
    \item \textbf{No sinais.} In traditional play, teammates use subtle facial signals to communicate hand strength. This channel is not available to LLM agents.
    \item \textbf{Prompt dependence.} LLM performance is sensitive to prompt format, as shown in our prompt sensitivity analysis. Results may vary with different prompt designs.
    \item \textbf{Temperature sensitivity.} Our default temperature of 0.7 balances exploration and exploitation, but optimal temperature may vary by model.
    \item \textbf{No learning.} LLMs play each game independently with no memory of prior games. Adaptive strategies over a tournament are not captured.
\end{itemize}
